\documentclass[11pt]{article}
\usepackage[utf8]{inputenc}
\usepackage{amsmath, amssymb, amsthm}
\usepackage{geometry}
\geometry{a4paper, margin=1in}
\usepackage{hyperref}

\title{Practice Problem Sets: Linearly Independent Sets and Bases}
\author{Generated Study Material}
\date{}

\begin{document}

\maketitle
\tableofcontents
\newpage

% ---------------------------------------------------------
\section{Concept Checklist \& Question Mapping}
% ---------------------------------------------------------
This section outlines the core concepts and problem archetypes extracted from the source material, along with the specific problems that test each concept.

\begin{itemize}
    \item \textbf{Concept 1: Verifying a Basis for $\mathbb{R}^n$} (Checking independence and span via pivots) \\
    \textit{Tested in Questions: 1, 2, 3}
    
    \item \textbf{Concept 2: Dimension Check \& Zero Vector Rule} (Identifying automatic linear dependence) \\
    \textit{Tested in Questions: 4, 5, 6}
    
    \item \textbf{Concept 3: Finding a Basis for Nul $A$} (Solving $A\mathbf{x} = \mathbf{0}$ and extracting parametric vectors) \\
    \textit{Tested in Questions: 7, 8, 9, 10}
    
    \item \textbf{Concept 4: Finding a Basis for Col $A$} (Using the "Original Column" Rule) \\
    \textit{Tested in Questions: 11, 12, 13, 14}
    
    \item \textbf{Concept 5: Finding a Basis for Row $A$} (Using non-zero rows of the echelon form) \\
    \textit{Tested in Questions: 15, 16}
    
    \item \textbf{Concept 6: Finding a Basis for a Spanned Space} (Translating span to column space) \\
    \textit{Tested in Questions: 17, 18}
    
    \item \textbf{Concept 7: Linear Independence in Polynomial Spaces ($\mathbb{P}_n$)} (Setting up systems from coefficients) \\
    \textit{Tested in Questions: 19, 20, 21}
    
    \item \textbf{Concept 8: Finding a Basis for a Polynomial Span} (Using the Spanning Set Theorem to discard redundant polynomials) \\
    \textit{Tested in Questions: 22, 23}
    
    \item \textbf{Concept 9: Extending a Linearly Independent Set to a Basis} (Appending standard basis vectors and row reducing) \\
    \textit{Tested in Questions: 24, 25, 26}
    
    \item \textbf{Concept 10: True/False Conceptual Analysis \& Pitfalls} (RREF Column Trap, Subspace Basis, Evaluating Polynomials at a Point) \\
    \textit{Tested in Questions: 27, 28, 29, 30}
\end{itemize}

\newpage

% ---------------------------------------------------------
\section{Problem Sets}
% ---------------------------------------------------------

\subsection*{Concept 1: Verifying a Basis for $\mathbb{R}^n$}
\textbf{1.} Determine if the set of vectors $S = \left\{ \begin{bmatrix} 2 \\ 5 \end{bmatrix}, \begin{bmatrix} -1 \\ 3 \end{bmatrix} \right\}$ forms a basis for $\mathbb{R}^2$. Justify your answer.

\vspace{0.5cm}
\textbf{2.} Determine if the set of vectors $S = \left\{ \begin{bmatrix} 1 \\ -1 \\ 0 \end{bmatrix}, \begin{bmatrix} 0 \\ 1 \\ -1 \end{bmatrix}, \begin{bmatrix} -1 \\ 0 \\ 1 \end{bmatrix} \right\}$ forms a basis for $\mathbb{R}^3$. Justify your answer.

\vspace{0.5cm}
\textbf{3.} Determine if the set of vectors $S = \left\{ \begin{bmatrix} 1 \\ 0 \\ 0 \end{bmatrix}, \begin{bmatrix} 2 \\ 3 \\ 0 \end{bmatrix}, \begin{bmatrix} 4 \\ 5 \\ 6 \end{bmatrix} \right\}$ forms a basis for $\mathbb{R}^3$. Justify your answer.

\subsection*{Concept 2: Dimension Check \& Zero Vector Rule}
\textbf{4.} Without performing row reduction, determine if the set $S = \left\{ \begin{bmatrix} 1 \\ 2 \end{bmatrix}, \begin{bmatrix} 3 \\ 4 \end{bmatrix}, \begin{bmatrix} 5 \\ 6 \end{bmatrix} \right\}$ is a basis for $\mathbb{R}^2$. Explain.

\vspace{0.5cm}
\textbf{5.} Let $S = \left\{ \begin{bmatrix} 0 \\ 0 \\ 0 \end{bmatrix}, \begin{bmatrix} 1 \\ 2 \\ 3 \end{bmatrix}, \begin{bmatrix} 4 \\ 5 \\ 6 \end{bmatrix} \right\}$. Is $S$ a basis for $\mathbb{R}^3$? Explain using a specific theorem.

\vspace{0.5cm}
\textbf{6.} You are given a set of 5 vectors in $\mathbb{R}^4$. Can this set be linearly independent? Can it be a basis for $\mathbb{R}^4$? Explain.

\subsection*{Concept 3: Finding a Basis for Nul $A$}
\textbf{7.} Find a basis for the null space of the matrix $A = \begin{bmatrix} 1 & 2 & 0 & -1 \\ 0 & 0 & 1 & 2 \end{bmatrix}$.

\vspace{0.5cm}
\textbf{8.} Find a basis for Nul $A$ where $A = \begin{bmatrix} 2 & -4 & 1 & 0 \\ -4 & 8 & -2 & 0 \end{bmatrix}$. (Hint: Scale your basis vectors to avoid fractions).

\vspace{0.5cm}
\textbf{9.} Find a basis for Nul $A$ where $A = \begin{bmatrix} 1 & 0 & 0 \\ 0 & 1 & 0 \\ 0 & 0 & 1 \end{bmatrix}$.

\vspace{0.5cm}
\textbf{10.} Find a basis for Nul $A$ where $A = \begin{bmatrix} 1 & -1 & 2 \\ 2 & -2 & 4 \end{bmatrix}$.

\subsection*{Concept 4: Finding a Basis for Col $A$}
\textbf{11.} Find a basis for the column space of $A = \begin{bmatrix} 1 & 2 & 3 \\ 4 & 5 & 6 \\ 7 & 8 & 9 \end{bmatrix}$.

\vspace{0.5cm}
\textbf{12.} Find a basis for Col $A$ where $A = \begin{bmatrix} 1 & -1 & 0 & 2 \\ 0 & 1 & -1 & 3 \\ 0 & 0 & 1 & 4 \end{bmatrix}$.

\vspace{0.5cm}
\textbf{13.} Find a basis for Col $A$ where $A = \begin{bmatrix} 2 & 4 \\ 3 & 6 \end{bmatrix}$.

\vspace{0.5cm}
\textbf{14.} Find a basis for Col $A$ where $A = \begin{bmatrix} 0 & 0 & 1 \\ 0 & 0 & 2 \\ 0 & 0 & 3 \end{bmatrix}$.

\subsection*{Concept 5: Finding a Basis for Row $A$}
\textbf{15.} Find a basis for the row space of $A = \begin{bmatrix} 1 & 2 & 3 \\ 4 & 5 & 6 \\ 7 & 8 & 9 \end{bmatrix}$.

\vspace{0.5cm}
\textbf{16.} Find a basis for Row $A$ where $A = \begin{bmatrix} 1 & -1 & 2 \\ 0 & 1 & -1 \\ 1 & 0 & 1 \end{bmatrix}$.

\subsection*{Concept 6: Finding a Basis for a Spanned Space}
\textbf{17.} Find a basis for the space spanned by the vectors $\mathbf{v}_1 = \begin{bmatrix} 1 \\ 1 \\ 1 \end{bmatrix}$, $\mathbf{v}_2 = \begin{bmatrix} 2 \\ 2 \\ 2 \end{bmatrix}$, and $\mathbf{v}_3 = \begin{bmatrix} 1 \\ 2 \\ 3 \end{bmatrix}$.

\vspace{0.5cm}
\textbf{18.} Find a basis for Span$\{\mathbf{v}_1, \mathbf{v}_2, \mathbf{v}_3\}$ where $\mathbf{v}_1 = \begin{bmatrix} 1 \\ 0 \\ -1 \end{bmatrix}$, $\mathbf{v}_2 = \begin{bmatrix} 0 \\ 1 \\ 1 \end{bmatrix}$, and $\mathbf{v}_3 = \begin{bmatrix} 1 \\ 1 \\ 0 \end{bmatrix}$.

\subsection*{Concept 7: Linear Independence in Polynomial Spaces}
\textbf{19.} Determine if the polynomials $\mathbf{p}_1(t) = 1+t$, $\mathbf{p}_2(t) = t+t^2$, and $\mathbf{p}_3(t) = 1+t^2$ form a linearly independent set in $\mathbb{P}_2$.

\vspace{0.5cm}
\textbf{20.} Are the polynomials $\mathbf{p}_1(t) = 2-t$ and $\mathbf{p}_2(t) = 4-2t$ linearly independent? Why or why not?

\vspace{0.5cm}
\textbf{21.} Consider $\mathbf{p}_1(t) = 1$, $\mathbf{p}_2(t) = 1+t$, and $\mathbf{p}_3(t) = 1+t+t^2$. Are they linearly independent in $\mathbb{P}_2$?

\subsection*{Concept 8: Finding a Basis for a Polynomial Span}
\textbf{22.} Find a basis for Span$\{\mathbf{p}_1, \mathbf{p}_2, \mathbf{p}_3\}$ where $\mathbf{p}_1(t) = 1+t$, $\mathbf{p}_2(t) = 1-t$, and $\mathbf{p}_3(t) = 4$.

\vspace{0.5cm}
\textbf{23.} Find a basis for the subspace of $\mathbb{P}_2$ spanned by $\{t^2, t^2+t, t^2+t+1, t\}$.

\subsection*{Concept 9: Extending a Linearly Independent Set to a Basis}
\textbf{24.} Let $H = \mathbb{R}^3$. You are given a linearly independent set $S = \left\{ \begin{bmatrix} 1 \\ 1 \\ 0 \end{bmatrix}, \begin{bmatrix} 0 \\ 1 \\ 1 \end{bmatrix} \right\}$. Extend $S$ to a basis for $\mathbb{R}^3$.

\vspace{0.5cm}
\textbf{25.} Extend the set $S = \left\{ \begin{bmatrix} 2 \\ 0 \\ 0 \\ 0 \end{bmatrix}, \begin{bmatrix} 0 \\ 3 \\ 0 \\ 0 \end{bmatrix} \right\}$ to a basis for $\mathbb{R}^4$.

\vspace{0.5cm}
\textbf{26.} Extend the single vector $\mathbf{v}_1 = \begin{bmatrix} 1 \\ -1 \\ 0 \end{bmatrix}$ to a basis for $\mathbb{R}^3$.

\subsection*{Concept 10: True/False Conceptual Analysis \& Pitfalls}
\textbf{27.} (T/F) A set of 4 vectors in $\mathbb{R}^5$ must be linearly independent because $p < n$. Explain.

\vspace{0.5cm}
\textbf{28.} (T/F) To find a basis for the column space of a matrix $A$, you row reduce $A$ to its RREF and take the pivot columns of the RREF matrix. Explain the flaw.

\vspace{0.5cm}
\textbf{29.} (T/F) If a set of vectors spans a subspace $H$, then it is automatically a basis for $H$. Explain.

\vspace{0.5cm}
\textbf{30.} (T/F) To check if polynomials $\mathbf{p}_1(t)$ and $\mathbf{p}_2(t)$ are linearly independent, we can evaluate them at $t=0$. If $c_1\mathbf{p}_1(0) + c_2\mathbf{p}_2(0) = 0$ implies $c_1=c_2=0$, then the polynomials are independent. Explain the flaw.

\newpage

% ---------------------------------------------------------
\section{Detailed Solutions}
% ---------------------------------------------------------

\textbf{Solution 1:}
To be a basis for $\mathbb{R}^2$, the two vectors must be linearly independent and span $\mathbb{R}^2$. We place them into a matrix $A$ and row reduce:
$$ A = \begin{bmatrix} 2 & -1 \\ 5 & 3 \end{bmatrix} $$
Calculate the determinant: $\det(A) = (2)(3) - (-1)(5) = 6 + 5 = 11$. Since the determinant is non-zero, the matrix is invertible. By the Invertible Matrix Theorem, its columns are linearly independent and span $\mathbb{R}^2$. 
\textbf{Final Answer:} Yes, the set is a basis for $\mathbb{R}^2$.

\vspace{0.5cm}
\textbf{Solution 2:}
Place the vectors into a matrix $A$:
$$ A = \begin{bmatrix} 1 & 0 & -1 \\ -1 & 1 & 0 \\ 0 & -1 & 1 \end{bmatrix} $$
Row reduce $A$:
$R_2 \leftarrow R_2 + R_1 \implies \begin{bmatrix} 1 & 0 & -1 \\ 0 & 1 & -1 \\ 0 & -1 & 1 \end{bmatrix}$
$R_3 \leftarrow R_3 + R_2 \implies \begin{bmatrix} 1 & 0 & -1 \\ 0 & 1 & -1 \\ 0 & 0 & 0 \end{bmatrix}$
There are only 2 pivots. The columns are linearly dependent and do not span $\mathbb{R}^3$. (Notice that $\mathbf{v}_1 + \mathbf{v}_2 + \mathbf{v}_3 = \mathbf{0}$).
\textbf{Final Answer:} No, the set is not a basis for $\mathbb{R}^3$.

\vspace{0.5cm}
\textbf{Solution 3:}
Place the vectors into a matrix $A$:
$$ A = \begin{bmatrix} 1 & 2 & 4 \\ 0 & 3 & 5 \\ 0 & 0 & 6 \end{bmatrix} $$
The matrix is already in upper triangular form. The diagonal entries (1, 3, 6) are all non-zero, meaning there is a pivot in every column and every row. Thus, the vectors are linearly independent and span $\mathbb{R}^3$.
\textbf{Final Answer:} Yes, the set is a basis for $\mathbb{R}^3$.

\vspace{0.5cm}
\textbf{Solution 4:}
We have a set of $p=3$ vectors in $\mathbb{R}^n$ where $n=2$. By the dimension theorem, any set containing more vectors than there are entries in each vector ($p > n$) is automatically linearly dependent. Since a basis must be linearly independent, this set cannot be a basis.
\textbf{Final Answer:} No, because $p > n$ (3 > 2), the set is linearly dependent.

\vspace{0.5cm}
\textbf{Solution 5:}
The set contains the zero vector $\mathbf{0}$. Any set containing the zero vector is automatically linearly dependent because a non-trivial solution to $c_1\mathbf{v}_1 + c_2\mathbf{v}_2 + c_3\mathbf{v}_3 = \mathbf{0}$ exists (e.g., $1\cdot\mathbf{0} + 0\cdot\mathbf{v}_2 + 0\cdot\mathbf{v}_3 = \mathbf{0}$).
\textbf{Final Answer:} No, the set is not a basis because it contains the zero vector, making it linearly dependent.

\vspace{0.5cm}
\textbf{Solution 6:}
We have $p=5$ vectors in $\mathbb{R}^4$ ($n=4$). Since $p > n$, the set is automatically linearly dependent. Because it is linearly dependent, it cannot be a basis for $\mathbb{R}^4$.
\textbf{Final Answer:} No, it cannot be linearly independent, and therefore cannot be a basis.

\vspace{0.5cm}
\textbf{Solution 7:}
The matrix $A$ is already in RREF. The pivot variables are $x_1$ and $x_3$. The free variables are $x_2$ and $x_4$.
Equations:
$x_1 + 2x_2 - x_4 = 0 \implies x_1 = -2x_2 + x_4$
$x_3 + 2x_4 = 0 \implies x_3 = -2x_4$
Parametric vector form:
$$ \mathbf{x} = x_2 \begin{bmatrix} -2 \\ 1 \\ 0 \\ 0 \end{bmatrix} + x_4 \begin{bmatrix} 1 \\ 0 \\ -2 \\ 1 \end{bmatrix} $$
\textbf{Final Answer:} Basis for Nul $A$ is $\left\{ \begin{bmatrix} -2 \\ 1 \\ 0 \\ 0 \end{bmatrix}, \begin{bmatrix} 1 \\ 0 \\ -2 \\ 1 \end{bmatrix} \right\}$.

\vspace{0.5cm}
\textbf{Solution 8:}
Row reduce $A = \begin{bmatrix} 2 & -4 & 1 & 0 \\ -4 & 8 & -2 & 0 \end{bmatrix}$.
$R_2 \leftarrow R_2 + 2R_1 \implies \begin{bmatrix} 2 & -4 & 1 & 0 \\ 0 & 0 & 0 & 0 \end{bmatrix}$.
Divide $R_1$ by 2: $\begin{bmatrix} 1 & -2 & 0.5 & 0 \\ 0 & 0 & 0 & 0 \end{bmatrix}$.
Pivot is $x_1$. Free variables are $x_2, x_3, x_4$.
$x_1 = 2x_2 - 0.5x_3$.
$$ \mathbf{x} = x_2 \begin{bmatrix} 2 \\ 1 \\ 0 \\ 0 \end{bmatrix} + x_3 \begin{bmatrix} -0.5 \\ 0 \\ 1 \\ 0 \end{bmatrix} + x_4 \begin{bmatrix} 0 \\ 0 \\ 0 \\ 1 \end{bmatrix} $$
To avoid fractions, scale the second vector by 2.
\textbf{Final Answer:} Basis for Nul $A$ is $\left\{ \begin{bmatrix} 2 \\ 1 \\ 0 \\ 0 \end{bmatrix}, \begin{bmatrix} -1 \\ 0 \\ 2 \\ 0 \end{bmatrix}, \begin{bmatrix} 0 \\ 0 \\ 0 \\ 1 \end{bmatrix} \right\}$.

\vspace{0.5cm}
\textbf{Solution 9:}
The matrix $A$ is the $3 \times 3$ identity matrix. The equation $A\mathbf{x} = \mathbf{0}$ has only the trivial solution $\mathbf{x} = \mathbf{0}$. The null space is the zero subspace $\{\mathbf{0}\}$. The zero subspace has no basis (or its basis is the empty set).
\textbf{Final Answer:} The basis is the empty set $\emptyset$.

\vspace{0.5cm}
\textbf{Solution 10:}
Row reduce $A = \begin{bmatrix} 1 & -1 & 2 \\ 2 & -2 & 4 \end{bmatrix}$.
$R_2 \leftarrow R_2 - 2R_1 \implies \begin{bmatrix} 1 & -1 & 2 \\ 0 & 0 & 0 \end{bmatrix}$.
Pivot is $x_1$. Free variables are $x_2, x_3$.
$x_1 = x_2 - 2x_3$.
$$ \mathbf{x} = x_2 \begin{bmatrix} 1 \\ 1 \\ 0 \end{bmatrix} + x_3 \begin{bmatrix} -2 \\ 0 \\ 1 \end{bmatrix} $$
\textbf{Final Answer:} Basis for Nul $A$ is $\left\{ \begin{bmatrix} 1 \\ 1 \\ 0 \end{bmatrix}, \begin{bmatrix} -2 \\ 0 \\ 1 \end{bmatrix} \right\}$.

\vspace{0.5cm}
\textbf{Solution 11:}
Row reduce $A = \begin{bmatrix} 1 & 2 & 3 \\ 4 & 5 & 6 \\ 7 & 8 & 9 \end{bmatrix}$.
$R_2 \leftarrow R_2 - 4R_1 \implies \begin{bmatrix} 1 & 2 & 3 \\ 0 & -3 & -6 \\ 7 & 8 & 9 \end{bmatrix}$
$R_3 \leftarrow R_3 - 7R_1 \implies \begin{bmatrix} 1 & 2 & 3 \\ 0 & -3 & -6 \\ 0 & -6 & -12 \end{bmatrix}$
$R_3 \leftarrow R_3 - 2R_2 \implies \begin{bmatrix} 1 & 2 & 3 \\ 0 & -3 & -6 \\ 0 & 0 & 0 \end{bmatrix}$
The pivots are in columns 1 and 2. We must select columns 1 and 2 from the \textit{original} matrix $A$.
\textbf{Final Answer:} Basis for Col $A$ is $\left\{ \begin{bmatrix} 1 \\ 4 \\ 7 \end{bmatrix}, \begin{bmatrix} 2 \\ 5 \\ 8 \end{bmatrix} \right\}$.

\vspace{0.5cm}
\textbf{Solution 12:}
The matrix $A = \begin{bmatrix} 1 & -1 & 0 & 2 \\ 0 & 1 & -1 & 3 \\ 0 & 0 & 1 & 4 \end{bmatrix}$ is already in echelon form. The pivots are clearly in columns 1, 2, and 3. We select columns 1, 2, and 3 from $A$.
\textbf{Final Answer:} Basis for Col $A$ is $\left\{ \begin{bmatrix} 1 \\ 0 \\ 0 \end{bmatrix}, \begin{bmatrix} -1 \\ 1 \\ 0 \end{bmatrix}, \begin{bmatrix} 0 \\ -1 \\ 1 \end{bmatrix} \right\}$.

\vspace{0.5cm}
\textbf{Solution 13:}
Row reduce $A = \begin{bmatrix} 2 & 4 \\ 3 & 6 \end{bmatrix}$.
$R_2 \leftarrow R_2 - \frac{3}{2}R_1 \implies \begin{bmatrix} 2 & 4 \\ 0 & 0 \end{bmatrix}$.
The pivot is in column 1. Select column 1 from the original matrix.
\textbf{Final Answer:} Basis for Col $A$ is $\left\{ \begin{bmatrix} 2 \\ 3 \end{bmatrix} \right\}$.

\vspace{0.5cm}
\textbf{Solution 14:}
The matrix $A = \begin{bmatrix} 0 & 0 & 1 \\ 0 & 0 & 2 \\ 0 & 0 & 3 \end{bmatrix}$ has only one non-zero column. Row reducing would yield a pivot in column 3.
\textbf{Final Answer:} Basis for Col $A$ is $\left\{ \begin{bmatrix} 1 \\ 2 \\ 3 \end{bmatrix} \right\}$.

\vspace{0.5cm}
\textbf{Solution 15:}
From Problem 11, the echelon form of $A$ is $\begin{bmatrix} 1 & 2 & 3 \\ 0 & -3 & -6 \\ 0 & 0 & 0 \end{bmatrix}$. The non-zero rows of the echelon form create a basis for the row space. (We can also scale the second row to $(0, 1, 2)$).
\textbf{Final Answer:} Basis for Row $A$ is $\{(1, 2, 3), (0, 1, 2)\}$.

\vspace{0.5cm}
\textbf{Solution 16:}
Row reduce $A = \begin{bmatrix} 1 & -1 & 2 \\ 0 & 1 & -1 \\ 1 & 0 & 1 \end{bmatrix}$.
$R_3 \leftarrow R_3 - R_1 \implies \begin{bmatrix} 1 & -1 & 2 \\ 0 & 1 & -1 \\ 0 & 1 & -1 \end{bmatrix}$.
$R_3 \leftarrow R_3 - R_2 \implies \begin{bmatrix} 1 & -1 & 2 \\ 0 & 1 & -1 \\ 0 & 0 & 0 \end{bmatrix}$.
The non-zero rows form the basis.
\textbf{Final Answer:} Basis for Row $A$ is $\{(1, -1, 2), (0, 1, -1)\}$.

\vspace{0.5cm}
\textbf{Solution 17:}
The space spanned by these vectors is the column space of $A = [\mathbf{v}_1 \ \mathbf{v}_2 \ \mathbf{v}_3]$.
$A = \begin{bmatrix} 1 & 2 & 1 \\ 1 & 2 & 2 \\ 1 & 2 & 3 \end{bmatrix}$.
Row reduce:
$R_2 \leftarrow R_2 - R_1 \implies \begin{bmatrix} 1 & 2 & 1 \\ 0 & 0 & 1 \\ 1 & 2 & 3 \end{bmatrix}$.
$R_3 \leftarrow R_3 - R_1 \implies \begin{bmatrix} 1 & 2 & 1 \\ 0 & 0 & 1 \\ 0 & 0 & 2 \end{bmatrix}$.
$R_3 \leftarrow R_3 - 2R_2 \implies \begin{bmatrix} 1 & 2 & 1 \\ 0 & 0 & 1 \\ 0 & 0 & 0 \end{bmatrix}$.
Pivots are in columns 1 and 3. Select $\mathbf{v}_1$ and $\mathbf{v}_3$.
\textbf{Final Answer:} Basis is $\left\{ \begin{bmatrix} 1 \\ 1 \\ 1 \end{bmatrix}, \begin{bmatrix} 1 \\ 2 \\ 3 \end{bmatrix} \right\}$.

\vspace{0.5cm}
\textbf{Solution 18:}
By inspection, notice that $\mathbf{v}_1 + \mathbf{v}_2 = \begin{bmatrix} 1 \\ 0 \\ -1 \end{bmatrix} + \begin{bmatrix} 0 \\ 1 \\ 1 \end{bmatrix} = \begin{bmatrix} 1 \\ 1 \\ 0 \end{bmatrix} = \mathbf{v}_3$.
Since $\mathbf{v}_3$ is a linear combination of $\mathbf{v}_1$ and $\mathbf{v}_2$, the Spanning Set Theorem allows us to discard it. $\mathbf{v}_1$ and $\mathbf{v}_2$ are not scalar multiples of each other, so they are independent.
\textbf{Final Answer:} Basis is $\left\{ \begin{bmatrix} 1 \\ 0 \\ -1 \end{bmatrix}, \begin{bmatrix} 0 \\ 1 \\ 1 \end{bmatrix} \right\}$.

\vspace{0.5cm}
\textbf{Solution 19:}
Set up the equation $c_1(1+t) + c_2(t+t^2) + c_3(1+t^2) = 0$.
Group by powers of $t$:
Constants: $c_1 + c_3 = 0$
$t$: $c_1 + c_2 = 0$
$t^2$: $c_2 + c_3 = 0$
Matrix form: $\begin{bmatrix} 1 & 0 & 1 \\ 1 & 1 & 0 \\ 0 & 1 & 1 \end{bmatrix}$.
Row reduce: $R_2 \leftarrow R_2 - R_1 \implies \begin{bmatrix} 1 & 0 & 1 \\ 0 & 1 & -1 \\ 0 & 1 & 1 \end{bmatrix}$.
$R_3 \leftarrow R_3 - R_2 \implies \begin{bmatrix} 1 & 0 & 1 \\ 0 & 1 & -1 \\ 0 & 0 & 2 \end{bmatrix}$.
There are 3 pivots. The only solution is the trivial solution.
\textbf{Final Answer:} Yes, they are linearly independent.

\vspace{0.5cm}
\textbf{Solution 20:}
Notice that $\mathbf{p}_2(t) = 4 - 2t = 2(2 - t) = 2\mathbf{p}_1(t)$. Since one polynomial is a scalar multiple of the other, the set is linearly dependent.
\textbf{Final Answer:} No, they are linearly dependent because $\mathbf{p}_2 = 2\mathbf{p}_1$.

\vspace{0.5cm}
\textbf{Solution 21:}
Set up the coordinate matrix relative to the standard basis $\{1, t, t^2\}$:
$A = \begin{bmatrix} 1 & 1 & 1 \\ 0 & 1 & 1 \\ 0 & 0 & 1 \end{bmatrix}$.
The matrix is already upper triangular with non-zero entries on the diagonal (3 pivots). Thus, the columns are linearly independent.
\textbf{Final Answer:} Yes, they are linearly independent.

\vspace{0.5cm}
\textbf{Solution 22:}
We check for linear dependence. Notice that $\mathbf{p}_1(t) + \mathbf{p}_2(t) = (1+t) + (1-t) = 2$.
Also, $\mathbf{p}_3(t) = 4 = 2(\mathbf{p}_1(t) + \mathbf{p}_2(t))$.
Since $\mathbf{p}_3$ is a linear combination of $\mathbf{p}_1$ and $\mathbf{p}_2$, we can discard it by the Spanning Set Theorem. $\mathbf{p}_1$ and $\mathbf{p}_2$ are not scalar multiples of each other, so they are independent.
\textbf{Final Answer:} Basis is $\{1+t, 1-t\}$.

\vspace{0.5cm}
\textbf{Solution 23:}
Let the polynomials be $\mathbf{p}_1 = t^2$, $\mathbf{p}_2 = t^2+t$, $\mathbf{p}_3 = t^2+t+1$, $\mathbf{p}_4 = t$.
Notice that $\mathbf{p}_2 - \mathbf{p}_1 = (t^2+t) - t^2 = t = \mathbf{p}_4$. So $\mathbf{p}_4$ is redundant.
Also, $\mathbf{p}_3 = \mathbf{p}_2 + 1$. The constant $1$ is not in the set, but we can see that $\{\mathbf{p}_1, \mathbf{p}_4, 1\}$ would span the same space. Let's map to coordinates in $\{1, t, t^2\}$:
$A = \begin{bmatrix} 0 & 0 & 1 & 0 \\ 0 & 1 & 1 & 1 \\ 1 & 1 & 1 & 0 \end{bmatrix}$.
Row reduce: Swap $R_1, R_3$: $\begin{bmatrix} 1 & 1 & 1 & 0 \\ 0 & 1 & 1 & 1 \\ 0 & 0 & 1 & 0 \end{bmatrix}$.
Pivots are in columns 1, 2, and 3. We keep the first three polynomials.
\textbf{Final Answer:} Basis is $\{t^2, t^2+t, t^2+t+1\}$.

\vspace{0.5cm}
\textbf{Solution 24:}
Append the standard basis vectors $\mathbf{e}_1, \mathbf{e}_2, \mathbf{e}_3$ to the set and row reduce to find the pivot columns.
$A = \begin{bmatrix} 1 & 0 & 1 & 0 & 0 \\ 1 & 1 & 0 & 1 & 0 \\ 0 & 1 & 0 & 0 & 1 \end{bmatrix}$.
$R_2 \leftarrow R_2 - R_1 \implies \begin{bmatrix} 1 & 0 & 1 & 0 & 0 \\ 0 & 1 & -1 & 1 & 0 \\ 0 & 1 & 0 & 0 & 1 \end{bmatrix}$.
$R_3 \leftarrow R_3 - R_2 \implies \begin{bmatrix} 1 & 0 & 1 & 0 & 0 \\ 0 & 1 & -1 & 1 & 0 \\ 0 & 0 & 1 & -1 & 1 \end{bmatrix}$.
The pivots are in columns 1, 2, and 3. We keep the first three columns of the original matrix.
\textbf{Final Answer:} Extended Basis is $\left\{ \begin{bmatrix} 1 \\ 1 \\ 0 \end{bmatrix}, \begin{bmatrix} 0 \\ 1 \\ 1 \end{bmatrix}, \begin{bmatrix} 1 \\ 0 \\ 0 \end{bmatrix} \right\}$.

\vspace{0.5cm}
\textbf{Solution 25:}
The given vectors are scalar multiples of $\mathbf{e}_1$ and $\mathbf{e}_2$. To form a basis for $\mathbb{R}^4$, we need 4 linearly independent vectors. We can simply add $\mathbf{e}_3$ and $\mathbf{e}_4$. The matrix formed by these 4 vectors will be diagonal with non-zero entries, guaranteeing independence.
\textbf{Final Answer:} Extended Basis is $\left\{ \begin{bmatrix} 2 \\ 0 \\ 0 \\ 0 \end{bmatrix}, \begin{bmatrix} 0 \\ 3 \\ 0 \\ 0 \end{bmatrix}, \begin{bmatrix} 0 \\ 0 \\ 1 \\ 0 \end{bmatrix}, \begin{bmatrix} 0 \\ 0 \\ 0 \\ 1 \end{bmatrix} \right\}$.

\vspace{0.5cm}
\textbf{Solution 26:}
Append $\mathbf{e}_1, \mathbf{e}_2, \mathbf{e}_3$ to $\mathbf{v}_1$:
$A = \begin{bmatrix} 1 & 1 & 0 & 0 \\ -1 & 0 & 1 & 0 \\ 0 & 0 & 0 & 1 \end{bmatrix}$.
$R_2 \leftarrow R_2 + R_1 \implies \begin{bmatrix} 1 & 1 & 0 & 0 \\ 0 & 1 & 1 & 0 \\ 0 & 0 & 0 & 1 \end{bmatrix}$.
The pivots are in columns 1, 2, and 4. We select columns 1, 2, and 4 from the original matrix.
\textbf{Final Answer:} Extended Basis is $\left\{ \begin{bmatrix} 1 \\ -1 \\ 0 \end{bmatrix}, \begin{bmatrix} 1 \\ 0 \\ 0 \end{bmatrix}, \begin{bmatrix} 0 \\ 0 \\ 1 \end{bmatrix} \right\}$.

\vspace{0.5cm}
\textbf{Solution 27:}
\textbf{False.} The theorem states that if $p > n$, the set is \textit{automatically dependent}. However, if $p \le n$ (here $4 \le 5$), the set \textit{can} be independent, but it is not \textit{guaranteed} to be. For example, the set could contain the zero vector, or two identical vectors, making it dependent despite having fewer vectors than dimensions.

\vspace{0.5cm}
\textbf{Solution 28:}
\textbf{False.} This is the "RREF Column Trap". Row operations preserve the row space and the null space, but they fundamentally alter the column space. You must use the RREF matrix only to identify \textit{where} the pivots are, and then you must extract those specific columns from the \textit{original} matrix $A$.

\vspace{0.5cm}
\textbf{Solution 29:}
\textbf{False.} A set can span a subspace $H$ without being a basis if it contains redundant vectors (i.e., it is linearly dependent). To be a basis, the set must span $H$ \textit{and} be linearly independent.

\vspace{0.5cm}
\textbf{Solution 30:}
\textbf{False.} Linear independence of functions requires that the equation $c_1\mathbf{p}_1(t) + c_2\mathbf{p}_2(t) = 0$ holds for \textit{all} $t$, not just a single specific value like $t=0$. Evaluating at a single point might falsely make independent polynomials look dependent (e.g., $\mathbf{p}_1(t)=t$ and $\mathbf{p}_2(t)=t^2$ both evaluate to 0 at $t=0$, but they are independent functions).

\newpage

% ---------------------------------------------------------
\section{Concept Checklist \& Question Mapping (Review)}
% ---------------------------------------------------------
\begin{itemize}
    \item \textbf{Concept 1: Verifying a Basis for $\mathbb{R}^n$} \\
    \textit{Tested in Questions: 1, 2, 3}
    
    \item \textbf{Concept 2: Dimension Check \& Zero Vector Rule} \\
    \textit{Tested in Questions: 4, 5, 6}
    
    \item \textbf{Concept 3: Finding a Basis for Nul $A$} \\
    \textit{Tested in Questions: 7, 8, 9, 10}
    
    \item \textbf{Concept 4: Finding a Basis for Col $A$} \\
    \textit{Tested in Questions: 11, 12, 13, 14}
    
    \item \textbf{Concept 5: Finding a Basis for Row $A$} \\
    \textit{Tested in Questions: 15, 16}
    
    \item \textbf{Concept 6: Finding a Basis for a Spanned Space} \\
    \textit{Tested in Questions: 17, 18}
    
    \item \textbf{Concept 7: Linear Independence in Polynomial Spaces ($\mathbb{P}_n$)} \\
    \textit{Tested in Questions: 19, 20, 21}
    
    \item \textbf{Concept 8: Finding a Basis for a Polynomial Span} \\
    \textit{Tested in Questions: 22, 23}
    
    \item \textbf{Concept 9: Extending a Linearly Independent Set to a Basis} \\
    \textit{Tested in Questions: 24, 25, 26}
    
    \item \textbf{Concept 10: True/False Conceptual Analysis \& Pitfalls} \\
    \textit{Tested in Questions: 27, 28, 29, 30}
\end{itemize}

\end{document}